%!TEX program = xelatex
%% ============================================================
%%  CIS 2026 — Congress on Intelligent Systems
%%  NIT Warangal x SCRS — Springer LNNS
%%  FINAL SUBMISSION DRAFT
%%
%%  Title: Morphology-Aware Telugu News Summarization:
%%         Comparing Extractive and Abstractive Approaches
%%         under Resource Constraints
%%
%%  System: Saaram ({\telugufont సారం})
%% ============================================================

\documentclass[runningheads]{llncs}

%% ── Packages ─────────────────────────────────────────────────
\usepackage{fontspec}
\usepackage{polyglossia}
\setmainlanguage{english}
\setotherlanguage{telugu}
\newfontfamily\telugufont{Noto Serif Telugu}[Script=Telugu]
\usepackage{amsmath,amssymb}
\usepackage{graphicx}
\usepackage{xcolor}
\usepackage{booktabs}
\usepackage{multirow}
\usepackage{microtype}
\usepackage{enumitem}
\usepackage[section]{placeins}
\usepackage{url}
\usepackage{cite}
\usepackage[hidelinks]{hyperref}

%% ── Float spacing ────────────────────────────────────────────
\setlength{\textfloatsep}{5pt}
\setlength{\floatsep}{4pt}
\setlength{\intextsep}{4pt}

%% ── TikZ ─────────────────────────────────────────────────────
\usepackage{tikz}
\usetikzlibrary{shapes.geometric,arrows.meta,positioning,
                fit,backgrounds,calc}
\tikzset{
  block/.style={rectangle,rounded corners=2pt,
    minimum width=2.0cm, minimum height=0.55cm,
    text centered, draw=black, fill=#1,
    font=\footnotesize\bfseries,
    text width=1.95cm, inner sep=2pt},
  block/.default=white,
  wideblock/.style={rectangle,rounded corners=2pt,
    minimum width=3.8cm, minimum height=0.70cm,
    text centered, draw=black, fill=#1,
    font=\small\bfseries,
    text width=3.6cm, inner sep=3pt},
  wideblock/.default=white,
  arrow/.style={-{Stealth[length=4pt]},thick},
  dasharrow/.style={-{Stealth[length=4pt]},thick,dashed},
  groupbox/.style={rectangle,rounded corners=3pt,
    draw=gray!70,dashed,thick,fill=#1!8,inner sep=5pt},
  groupbox/.default=gray
}

%% ============================================================
\begin{document}
%% ============================================================

\title{Saaram: Resource-Aware Telugu News Summarization with Morphology-Aware TF-IDF and mT5}

\titlerunning{Saaram: Resource-Aware Telugu News Summarization}
\authorrunning{Hariharan et al.}

\author{
  Hariharan Narlakanti\inst{1} \and
  Vivek Nidumolu\inst{1}       \and
  Vishnu Vardhan Reddy Padala\inst{1} \and
  Sanjeev Practur\inst{1}      \and
  Vinay Simha Reddy Tappeta\inst{1}
}

\institute{
  Department of Computer Science and Engineering (AI \& ML),\\
  Malla Reddy University, Hyderabad, Telangana, India\\
  \email{narlakantihariharan@gmail.com}
}

\maketitle

%% ============================================================
%%  ABSTRACT
%% ============================================================
\begin{abstract}
Telugu is a morphologically rich, agglutinative language that poses significant challenges for tokenization and evaluation in standard NLP pipelines. We present
\textit{Saaram} ({\telugufont సారం}, ``essence''), a deployed Telugu news summarization service with optional speech synthesis, and conduct a comparative study of five configurations evaluated on the complete 1,302-sample XL-Sum Telugu test split. Our morphology-aware TF-IDF baseline replaces word-level tokenization with character $n$-gram analysis (\texttt{analyzer='char\_wb'}, $n \in [2,4]$) and lead-bias position weighting, achieving a 26.6\% relative improvement in ROUGE-2 over the word-level baseline. Evaluated abstractive models show that the 8-bit LoRA fine-tuned mT5 model achieves the highest lexical overlap (ROUGE-2: 0.0496) but yields near-identical semantic quality, measured by BERTScore (0.7272), compared with mT5 Base (0.7273), indicating that in-distribution fine-tuning provides negligible semantic gains. Finally, a resource-aware runtime router adaptively dispatches requests between TF-IDF and mT5 Base, retaining 99.56\% of mT5 semantic quality (BERTScore) while mitigating memory exhaustion risk. Latency benchmarking further characterizes deployment-relevant cost--quality trade-offs under resource constraints.
\end{abstract}

\keywords{Telugu text summarization \and morphology-aware TF-IDF \and mT5 \and low-resource NLP \and resource-aware routing}

%% ============================================================
\section{Introduction}
%% ============================================================

Telugu, with over 80~million native speakers across Telangana
and Andhra Pradesh, is one of the most widely spoken languages
in India, yet remains severely underrepresented in deployed NLP systems~\cite{joshi2020text}. The rapid proliferation of
digital Telugu news creates a clear societal need: readers face an overwhelming volume of content, while few publicly accessible automated summarization systems exist for Telugu at scale.

Two structural challenges compound this problem. First, Telugu
is an \emph{agglutinative} language, where a single root word can generate dozens of inflected surface forms through affixation. Standard NLP
pipelines that assume word-level tokenization—primarily designed for
morphologically simpler languages such as English—treat these
inflected forms as unrelated vocabulary items, severely
underestimating term importance and fragmenting evaluation
metrics~\cite{singh2020text}. For example, the root
\textit{{\telugufont రాష్ట్రం}} (state) appears in news text as
\textit{{\telugufont రాష్ట్రంలో}} (in the state) and
\textit{{\telugufont రాష్ట్రానికి}} (to the state), yet
word-level TF-IDF treats all three as entirely unrelated tokens.
Second, many standard ROUGE implementations strip non-ASCII characters by default, rendering Telugu ROUGE evaluation invalid without explicit adaptation.

Prior Telugu NLP research has addressed these challenges in
isolation: extractive approaches without morphological
awareness~\cite{anandbabu2023extractive}, or neural
abstractive models without deployment-oriented
analysis~\cite{srinivasa2019neural}. To the best of our knowledge,
no prior Telugu summarization study jointly evaluates
morphology-aware extractive methods, large-scale neural
summarization, and production deployment trade-offs within a unified framework.

This paper presents \textit{Saaram} ({\telugufont సారం}, ``essence''), a
deployed Telugu news summarization system with optional speech synthesis, and makes five contributions:

\begin{itemize}[leftmargin=*,itemsep=1pt,topsep=2pt]
  \item A \textbf{morphology-aware TF-IDF baseline} (TF-IDF~V2) using character $n$-grams and lead-bias weighting, achieving a 26.6\% relative improvement in ROUGE-2 over word-level TF-IDF~V1.
  
  \item A \textbf{large-scale comparative benchmark} evaluating five summarization configurations on the complete 1,302-sample XL-Sum Telugu test split using ROUGE, BERTScore, and latency profiling.
  
  \item A \textbf{LoRA fine-tuning ablation study} showing that 8-bit quantized LoRA adaptation yields only marginal lexical gains (+0.0032 ROUGE-L) without semantic improvement (BERTScore: 0.7272 vs.\ 0.7273) over the pre-trained base model.
  
  \item A \textbf{resource-aware inference strategy} that routes requests at runtime while retaining 99.56\% of mT5 semantic quality (BERTScore), routing only 5.38\% of samples (70 of 1,302) to TF-IDF~V2 as a robust fallback under resource constraints.
  
  \item A \textbf{resource-constrained deployment analysis} on free-tier CPU cloud infrastructure, characterizing memory-aware lazy loading, fallback behavior, and TTS latency.
\end{itemize}

The complete system is publicly deployed, with the frontend hosted on Vercel, the backend on Hugging Face Spaces (Docker, CPU-only), and the source code available on GitHub.
%% ============================================================
\section{Related Work}
%% ============================================================

\textbf{Extractive and neural abstractive summarization.}
Graph-based methods such as LexRank~\cite{erkan2004lexrank} and TextRank~\cite{mihalcea2004textrank} rank sentence salience using centrality measures, while MMR~\cite{carbonell1998use} introduces diversity-aware reranking. TF-IDF-based sentence scoring remains a classical baseline for both English and Indic summarization tasks~\cite{joshi2020text}. Prior Telugu extractive work~\cite{anandbabu2023extractive} primarily relies on word-level tokenization, which does not adequately capture agglutinative morphology. In abstractive summarization, attention-based architectures~\cite{rush2015neural}, pointer-copy mechanisms~\cite{nallapati2016abstractive}, and coverage penalties~\cite{see2017get} have improved sequence-to-sequence generation. The Transformer~\cite{vaswani2017attention} and T5~\cite{raffel2020exploring} further enabled unified text-to-text modeling. Multilingual T5 (mT5)~\cite{xue2021mt5} extends this framework to 101 languages, using SentencePiece-based subword segmentation that better accommodates Telugu morphology.


\textbf{Indic and Telugu NLP.}
IndicBART~\cite{kumar2022indicbart} is a pretrained model for Indic text generation, although our preliminary qualitative screening revealed instability on Telugu summarization inputs (Sect.~\ref{sec:model-selection}). Doddapaneni et al.~\cite{doddapaneni2021primer} survey Indic language models and identify data scarcity as a primary bottleneck, which datasets such as XL-Sum~\cite{hasan2021xlsum} and TeSum~\cite{urlana2022tesum} help alleviate. Prior Telugu summarization work~\cite{anandbabu2023extractive} applied TextRank and MMR while retaining word-level tokenization, a limitation addressed in this study.


\textbf{Evaluation for agglutinative languages.}
ROUGE~\cite{lin2004rouge} measures exact $n$-gram overlap, often penalizing inflected paraphrases common in agglutinative languages. BERTScore~\cite{zhang2020bertscore} uses pretrained contextual embeddings to compute token similarity, theoretically assigning partial credit to semantically equivalent yet morphologically varied tokens. Motivated by this property, we use BERTScore as the primary semantic similarity metric while reporting ROUGE as a supplementary lexical overlap measure.


\textbf{Deployment gap.}
Indic NLP literature rarely reports real-world deployment characteristics. Existing systems typically emphasize quality metrics without analyzing runtime latency, memory constraints, or resource-aware fallback mechanisms. This gap is addressed in our work through comprehensive deployment-aware profiling.


%% ============================================================
\section{System Architecture}
\label{sec:arch}
%% ============================================================

\textit{Saaram} implements a three-tier architecture designed under two practical deployment constraints: free-tier cloud hosting with limited memory and production CPU-only serving constraints, together with interactive latency requirements for user-facing summarization. Fig.~\ref{fig:arch} illustrates the complete pipeline. The frontend exposes three inference quality tiers corresponding to TF-IDF~V2, mT5~Base, and mT5~Fine-tuned models. Optional Edge TTS audio synthesis is supported for spoken summaries.


\begin{figure}[t]
\centering
\resizebox{\columnwidth}{!}{%
\begin{tikzpicture}[node distance=0.9cm and 0.7cm]

  %% Tier 1: Frontend
  \node[wideblock=blue!20,minimum height=0.9cm] (fe)
    {\textit{Saaram} Frontend\\
     {\small React.js $|$ Method Select $|$ Audio}};

  %% Tier 2: Backend
  \node[wideblock=orange!20,minimum height=0.9cm,
        right=1.8cm of fe] (api)
    {FastAPI Backend\\
     {\small Pydantic $|$ Router $|$ Dispatch}};

  \draw[arrow](fe.east)--(api.west)
    node[midway,above,font=\scriptsize\itshape]{JSON/Axios};

  %% Tier 3: Pipeline nodes
  \node[block=purple!15,text width=2.0cm,
        below=1.9cm of fe] (ext)
    {Extraction\\{\scriptsize BS4 / RSS}};

  \node[block=purple!15,text width=2.0cm,
        right=0.5cm of ext] (cln)
    {Unicode\\Cleaning};

  \node[block=yellow!25,text width=2.0cm,
        right=0.5cm of cln] (rtr)
    {Adaptive\\Router};

  \node[block=green!15,text width=2.0cm,
        right=0.5cm of rtr] (tfi)
    {TF-IDF V2\\{\scriptsize char n-gram}};

  \node[block=green!25,text width=2.0cm,
        right=0.5cm of tfi] (mb)
    {mT5 Base\\{\scriptsize XLSum}};

  \node[block=green!35,text width=2.0cm,
        right=0.5cm of mb] (mft)
    {mT5 Fine\\Tuned};

  \node[block=cyan!20,text width=2.2cm,
        right=0.5cm of mft] (tts)
    {Edge TTS\\{\scriptsize te-IN-Shruti}};

  %% Connections
  \draw[arrow](ext)--(cln);
  \draw[arrow](cln)--(rtr);
  \draw[arrow](rtr)--(tfi);
  \draw[arrow](rtr.north)--++(0,0.4)-|(mb.north);
  \draw[dasharrow](mb)--(mft)
    node[midway,below,font=\tiny\itshape]{optional};
  \draw[arrow](tfi.south)--++(0,-0.3)-|(tts.south);
  \draw[arrow](mb.south)--++(0,-0.3)-|(tts.south);

  \draw[arrow](api.south)--++(0,-0.65)-|(cln.north);
  \draw[arrow](tts.north)--++(0,0.9)-|(api.south);

  %% Group boxes
  \begin{scope}[on background layer]
    \node[groupbox=purple,fit=(ext)(cln),
          inner sep=5pt](gi){};
    \node[groupbox=yellow,fit=(rtr),
          inner sep=5pt](gr){};
    \node[groupbox=green,fit=(tfi)(mb)(mft),
          inner sep=5pt](gs){};
    \node[groupbox=cyan,fit=(tts),
          inner sep=6pt](gt){};
  \end{scope}

  \node[below=3pt of gi,font=\scriptsize\itshape]
    {Input Processing};
  \node[below=3pt of gr,font=\scriptsize\itshape]
    {Router};
  \node[below=3pt of gs,
        font=\tiny\itshape,text width=3.6cm,align=center]
    {Summarization (3 methods)};
  \node[below=3pt of gt,font=\scriptsize\itshape]
    {TTS};

\end{tikzpicture}}
\caption{\textit{Saaram} three-tier architecture. The resource-aware
router selects TF-IDF~V2 or mT5~Base at runtime; mT5~Fine-tuned
(8-bit LoRA) is available as a manual override. Edge TTS converts summaries
to Telugu audio asynchronously.}
\label{fig:arch}
\end{figure}

\subsection{Input Processing}
\label{sec:input}

The system supports three input modes: (1)~direct text entry, (2)~URL-based article extraction using BeautifulSoup4 with a 10-second timeout and \texttt{<p>}-element prioritization, and (3)~live RSS ingestion from BBC Telugu for the Saaram Radio bulletin feature. All inputs pass through a Unicode-aware regular-expression cleaning pipeline that preserves the Telugu Unicode block (U+0C00--U+0C7F), dandas (U+0964), ASCII digits, and standard punctuation. This preprocessing pipeline is applied identically during fine-tuning and inference to ensure training--inference consistency.


\subsection{Morphology-Aware TF-IDF (V2)}
\label{sec:tfidf}

Standard word-level TF-IDF (V1) fails for Telugu because
agglutinative inflection creates distinct surface tokens for
semantically related forms. Consider the root
\textit{{\telugufont రాష్ట్రం}} (state): its inflected forms
\textit{{\telugufont రాష్ట్రంలో}} (in the state) and
\textit{{\telugufont రాష్ట్రానికి}} (to the state) share no whitespace-
delimited word token, causing TF-IDF to assign separate, low
term weights to each form of a high-frequency concept.

TF-IDF~V2 addresses this with four targeted improvements:

(1)~\textbf{Character n-gram analysis:} replacing word tokenization with character n-grams (\texttt{analyzer='char\_wb'}, $\texttt{ngram\_range}=(2,4)$) to capture root substrings across inflected forms, restoring term-frequency coherence; (2)~\textbf{Lead-bias position weighting:} applying decaying position weights to scores to boost earlier sentences, reflecting news inverted-pyramid structures; (3)~\textbf{Adaptive sentence selection:} dynamically adjusting the sentence extraction count $k$ based on input length to prevent under/over-extraction; and (4)~\textbf{Per-sentence score exposure:} exposing sentence-level scores to the router (Sect.~\ref{sec:router}) as confidence signals.

\subsection{mT5 Abstractive Summarization}
\label{sec:mt5}

\textbf{mT5 Base.}
We use the
\texttt{csebuetnlp/mT5\_multilingual\_XLSum}
checkpoint (580M parameters), pretrained on multilingual
mC4 and further adapted for multilingual summarization on
the XL-Sum benchmark, including Telugu.
SentencePiece tokenization naturally captures Telugu
subword morphology.
Inference uses beam search with width~4, length
penalty~2.0, trigram blocking, output length bounds of
30--128 tokens, and a maximum input length of 512 tokens.

\textbf{mT5 Fine-tuned (8-bit LoRA).}
The base checkpoint is further adapted to XL-Sum Telugu
using parameter-efficient LoRA~\cite{hu2022lora}
applied to an 8-bit quantized model~\cite{dettmers2022int8}.
LoRA uses rank $r=16$, scaling factor $\alpha=32$,
target modules \texttt{[q,v]}, and dropout 0.05.

Training was performed for 3 epochs with learning rate
$3\times10^{-4}$ and per-device batch size 4 with
gradient accumulation $\times 2$ (effective batch size 8)
on a single NVIDIA Tesla T4 GPU (16\,GB VRAM, Kaggle
environment). Training input and target lengths were
capped at 384 and 96 tokens, respectively, while deployed
inference uses \texttt{max\_length=128}. The dataset was
split into 10,421 training, 1,302 validation, and 1,302
test samples. An earlier full-parameter fine-tuned
baseline (3.372\,s latency) is reported in
Table~\ref{tab:latency} for reference.

Since the base checkpoint already contains Telugu
summarization knowledge from XL-Sum pre-adaptation,
additional LoRA fine-tuning yields only marginal lexical
improvements (ROUGE-L +0.0032) without measurable
semantic gains (BERTScore: 0.7272 vs.\ 0.7273).


\subsection{Model Backbone Selection}
\label{sec:model-selection}

Four abstractive backbones were screened qualitatively
using a representative set of Telugu news articles,
evaluated by two native Telugu speakers and one NLP
researcher.

\textbf{IndicBART} (244M parameters) produced unstable
outputs, frequently exhibiting token repetition,
language switching, and named-entity hallucinations.
\textbf{mBART-50} (611M parameters) generated
grammatically valid summaries but often produced generic
or less informative outputs. \textbf{mT5 Base}
(580M parameters) produced fluent, language-consistent
summaries with strong semantic preservation.
\textbf{mT5 Fine-tuned} (580M parameters) showed similar
fluency while providing slightly improved domain
consistency.

Based on this qualitative screening, the two mT5
variants were selected for full quantitative evaluation.

\subsection{Resource-Aware Inference Strategy}
\label{sec:router}

As a deployment optimization layer, the system includes a
lightweight inference router that intercepts each
summarization request before model dispatch and selects
between TF-IDF~V2 and mT5~Base without user intervention.

Routing decisions are based on three signals:
(1)~\textbf{Word count}, where inputs below 75 words are
routed to TF-IDF~V2 and inputs above 200 words to
mT5~Base;
(2)~\textbf{TF-IDF score variance}, where variance above
0.04 favors TF-IDF~V2 and variance below 0.008 favors
mT5~Base; and
(3)~\textbf{Available RAM}, where low-memory conditions
force TF-IDF~V2 regardless of other signals.

The router never selects the fine-tuned model
automatically; fine-tuned inference remains a manual
override reserved for cases where marginal ROUGE gains
are desired.

On the XL-Sum test split, 70 of 1,302 samples (5.38\%)
were routed to TF-IDF~V2, reflecting the dataset bias
toward longer news articles. Since the router activates
for only a small fraction of in-distribution samples, it
primarily serves as a deployment safeguard for graceful
degradation under resource constraints rather than as a
core algorithmic contribution.


\subsection{Deployment Constraints and Optimization}
\label{sec:deployment}

\textit{Saaram} is deployed on Hugging Face Spaces under
free-tier constraints (2~vCPU, 16~GB RAM, CPU-only
inference), imposing strict resource limitations that
required explicit architectural optimizations.

\textbf{Memory-aware lazy loading.}
The mT5 checkpoints require approximately 2.3~GB RAM for
the 8-bit LoRA checkpoint and 3.7~GB for the earlier
full-parameter baseline. Loading multiple checkpoints
simultaneously alongside FastAPI, the TF-IDF pipeline,
and operating-system overhead risks out-of-memory (OOM)
failures. To mitigate this, models are initialized only
upon first request (\emph{lazy loading}) and garbage-
collected after a configurable idle timeout. As a result,
only the actively requested model occupies memory at any
given time.

\textbf{Robust fallback hierarchy.}
The system handles three major failure modes through
automatic degradation. First, OOM events during mT5
inference trigger immediate fallback to TF-IDF~V2, with
the API returning a flag indicating degraded mode.
Second, checkpoint loading failures trigger fallback from
the fine-tuned model to mT5 Base, and subsequently to
TF-IDF~V2 if required. Third, URL extraction failures
return structured error responses with clear recovery
guidance. These safeguards ensure that the system never
returns an unhandled empty response.

\textbf{API design.}
The backend exposes seven FastAPI endpoints across Health,
News, Summarization, and Audio services. Request schemas
use Pydantic validation, and invalid requests return
structured HTTP~422 error responses for improved
debuggability and client-side handling.

%% ============================================================
\section{Experimental Setup}
\label{sec:setup}
%% ============================================================

\textbf{Dataset.}
All experiments use the complete XL-Sum Telugu test
split~\cite{hasan2021xlsum}, consisting of 1,302 BBC
Telugu news articles paired with human-written reference
summaries across politics, business, sports, and current
affairs. The base mT5 checkpoint was pretrained on
multilingual mC4 and subsequently fine-tuned on the
multilingual XL-Sum corpus, including Telugu, making
this evaluation largely in-distribution (discussed in
Sect.~\ref{sec:analysis}).

\textbf{Evaluation pipeline.}
The evaluation pipeline processes the full 1,302-sample
test split with checkpointing/resume support, structured
CSV outputs, and mean/P95/P99 latency profiling.

\textbf{ROUGE adaptation.}
Standard ROUGE scorers strip non-ASCII characters,
making Telugu evaluation invalid. We therefore implement
a custom Unicode-preserving tokenizer for reliable
ROUGE computation over Telugu text.

\textbf{Preprocessing consistency.}
All configurations receive identical cleaned inputs using
the preprocessing pipeline described in
Sect.~\ref{sec:input}. The same preprocessing settings
are applied during evaluation and baseline scoring.

\textbf{Metrics.}
We report ROUGE-1, ROUGE-2, ROUGE-L, and BERTScore F1.
BERTScore is computed using
\texttt{bert-base-multilingual-cased} with default
rescaling for Telugu. Wall-clock latency, measured from
request ingestion to output serialization, is averaged
across the test split.

%% ============================================================
\section{Results}
\label{sec:results}
%% ============================================================

\subsection{TF-IDF V1 vs V2: Morphological Improvement}

We isolate the contribution of morphology-aware
character $n$-gram analysis by comparing TF-IDF~V1
(word-level) and TF-IDF~V2 (character $n$-gram) in
Tables~\ref{tab:main} and~\ref{tab:latency}. TF-IDF~V2
consistently improves all ROUGE metrics while preserving
near-zero inference cost.

ROUGE-2 shows the largest relative gain, improving by
26.6\% (0.0128 $\rightarrow$ 0.0162), suggesting that
character $n$-grams mitigate bigram fragmentation caused
by Telugu morphological inflection. ROUGE-1 and
ROUGE-L improve by 14.4\% (0.0709 $\rightarrow$ 0.0811)
and 13.4\% (0.0554 $\rightarrow$ 0.0628), respectively.
Despite these gains, mean latency remains effectively
negligible at 0.005\,s.

\subsection{Full Comparative Evaluation}

Table~\ref{tab:main} summarizes quality metrics across all
five evaluated configurations, enabling direct comparison
of extractive, abstractive, and hybrid approaches.
Table~\ref{tab:latency} reports inference latency
separately to preserve focus on summarization quality in
the primary comparison.

\begin{table}[!htbp]
\centering
\caption{Quality Evaluation --- XL-Sum Telugu Test Split
(1,302 samples). Bold indicates best per metric.}
\label{tab:main}
\scriptsize
\setlength{\tabcolsep}{3.5pt}
\begin{tabular}{p{3.5cm}cccc}
\toprule
\textbf{Method}
& \textbf{ROUGE-1}
& \textbf{ROUGE-2}
& \textbf{ROUGE-L}
& \textbf{BERTScore} \\
\midrule
TF-IDF V1 (word-level)
& 0.0709 & 0.0128 & 0.0554 & 0.6626 \\

TF-IDF V2 (char $n$-gram)
& 0.0811 & 0.0162 & 0.0628 & 0.6671 \\

mT5 Base (pretrained)
& 0.1610 & 0.0478 & 0.1420 & \textbf{0.7273} \\

mT5 Fine-tuned$^\dagger$
& \textbf{0.1639}
& \textbf{0.0496}
& \textbf{0.1452}
& 0.7272 \\

Auto Router (hybrid)
& 0.1562 & 0.0458 & 0.1372 & 0.7241 \\
\midrule
\multicolumn{5}{p{11.2cm}}{\footnotesize
ROUGE values are reported as F1 scores.
BERTScore F1 uses \texttt{bert-base-multilingual-cased} with rescaling.
$^\dagger$ Fine-tuned model uses 8-bit LoRA adaptation.}\\
\bottomrule
\end{tabular}
\end{table}

\begin{table}[!htbp]
\centering
\caption{Inference Latency Profile --- Local MPS-accelerated evaluation (1,302 samples).}
\label{tab:latency}
\scriptsize
\setlength{\tabcolsep}{6pt}
\begin{tabular}{p{4.2cm}c}
\toprule
\textbf{Method} & \textbf{Mean Latency (s)} \\
\midrule
TF-IDF V1                & 0.001 \\
TF-IDF V2                & 0.005 \\
mT5 Base                 & 2.583 \\
mT5 Fine-tuned$^\dagger$ & 2.482 \\
Auto Router              & 3.130 \\
\midrule
\multicolumn{2}{p{8.3cm}}{\footnotesize
$^\dagger$ Fine-tuned model uses 8-bit LoRA adaptation.
Latency reduced versus the earlier full-parameter fine-tuned
checkpoint (3.372\,s; Sect.~\ref{sec:mt5}).} \\
\bottomrule
\end{tabular}
\end{table}

\textbf{(F1) mT5 significantly outperforms TF-IDF.}
mT5~Base achieves ROUGE-L of 0.1420 compared with
0.0628 for TF-IDF~V2, representing a 126\% relative
improvement and aligning with prior multilingual
summarization baselines~\cite{hasan2021xlsum}.
BERTScore similarly improves from 0.6671 to 0.7273.
TF-IDF~V2 consistently outperforms TF-IDF~V1,
confirming that morphology-aware tokenization improves
extractive summarization quality.

\textbf{(F2) LoRA fine-tuning provides negligible semantic gain.}
mT5 Fine-tuned (8-bit LoRA) yields near-identical
semantic quality to mT5~Base (BERTScore: 0.7272 vs.\
0.7273), indicating that additional adaptation on an
already covered distribution primarily affects lexical
realization rather than semantic representation.

\textbf{(F3) Inference latency remains practical.}
The 8-bit LoRA fine-tuned model achieves mean latency
of 2.482\,s, slightly faster than mT5~Base (2.583\,s)
and substantially faster than the earlier full-parameter
fine-tuned baseline (3.372\,s).

\textbf{(F4) Resource-aware routing retains 99.56\% of mT5 quality.}
The router preserves 99.56\% of mT5 quality
(BERTScore: 0.7241 vs.\ 0.7273) while routing only
5.38\% of samples to TF-IDF~V2, serving as an effective
fallback mechanism under resource-constrained deployment.



%% ============================================================
%% ============================================================
\section{Discussion and Analysis}
\label{sec:analysis}
%% ============================================================

\textbf{In-Distribution Fine-Tuning.}
The near-identical performance of mT5~Base and mT5~Fine-tuned (8-bit LoRA) indicates that local LoRA adaptation yields negligible semantic improvement over the base model. Since the base checkpoint (\texttt{csebuetnlp/mT5\_multilingual\_XLSum}) was already fine-tuned on the multilingual XL-Sum corpus, including Telugu, further in-distribution adaptation primarily alters surface phrasing rather than underlying semantic representations. Practitioners may therefore benefit more from fine-tuning on out-of-distribution domains where pretrained representations are weaker.

\textbf{Morphology-Aware Character N-Grams.}
The 26.6\% ROUGE-2 improvement of TF-IDF~V2 over TF-IDF~V1 highlights the value of morphology-aware tokenization. In Telugu, a concept such as \textit{government} can appear as inflected surface forms (\textit{{\telugufont ప్రభుత్వం}}, \textit{{\telugufont ప్రభుత్వానికి}}, \textit{{\telugufont ప్రభుత్వంపై}}) that share the root sequence \textit{{\telugufont ప్రభుత్వ}}. Character $n$-grams with range $(2,4)$ capture these shared substrings, aggregating term-frequency signals that word-level TF-IDF would otherwise fragment across distinct vocabulary entries, thereby improving salience estimation.

\textbf{Router Optimization.}
The inference router functions primarily as a deployment safeguard, routing only 5.38\% of XL-Sum test samples (70 of 1,302) to TF-IDF~V2. Because most in-distribution news articles are sufficiently long and information-dense, mT5~Base remains the dominant inference path. The router therefore serves mainly as an engineering fallback that improves robustness by enabling graceful degradation from neural to statistical summarization under resource constraints.

\textbf{BERTScore Suitability.}
BERTScore computes contextual token similarity and can assign partial credit to paraphrases that preserve meaning despite lexical variation. This property is especially relevant for morphologically agglutinative languages such as Telugu, where the same concept may appear in multiple inflected forms. We therefore use BERTScore as the primary semantic similarity metric, while acknowledging that its practical suitability relative to ROUGE for Telugu remains unverified without human evaluation. Such validation is an important direction for future work.

%% ============================================================
\section{Qualitative Analysis}
\label{sec:qualitative}
%% ============================================================
Table~\ref{tab:qualitative} presents a qualitative comparison
on a representative sample from the XL-Sum Telugu test split
(sample ID: \texttt{international-53333964}). System outputs
were generated using the local \textit{Saaram} evaluation
pipeline under the same preprocessing and inference settings
used for quantitative evaluation. The table shows top-ranked
extractive sentences for TF-IDF variants and decoded outputs
for mT5 variants, illustrating the qualitative difference
between extraction and abstraction.

\begin{table}[!htbp]
\centering
\caption{Qualitative Output Comparison --- XL-Sum Telugu test sample (index 0, \texttt{international-53333964}). Outputs are actual system inference results.}
\label{tab:qualitative}
\scriptsize
\setlength{\tabcolsep}{3pt}
\begin{tabular}{p{2.2cm}p{6.6cm}}
\toprule
\textbf{Source} & \textbf{Telugu Script} \\
\midrule

Input article &
{\telugufont ఎఫ్‌బీఐ డైరెక్టర్ క్రిస్టఫర్ రే “చైనా ప్రభుత్వ గూఢచర్యం, డేటా చోరీ వల్ల అమెరికా భవిష్యత్తుకు ఎప్పుడూ లేనంత దీర్ఘకాలిక ముప్పు ఉంది” అని అన్నారు. “చైనా చాలా స్థాయిల్లో ఆపరేషన్లు నిర్వహిస్తోంది.”}
\textit{[Article continues; full text available in XL-Sum Telugu dataset, sample ID: \texttt{international-53333964}.]} \\
\midrule

Gold reference &
{\telugufont చైనా ప్రభుత్వం నుంచి అమెరికాకు అతిపెద్ద దీర్ఘకాలిక ముప్పు పొంచి ఉందని అమెరికా నిఘా ఏజెన్సీ ఎఫ్‌బీఐ డైరెక్టర్ క్రిస్టఫర్ రే అన్నారు.} \\
\midrule

TF-IDF V1 &
{\telugufont ఎఫ్‌బీఐ డైరెక్టర్ క్రిస్టఫర్ రే “చైనా ప్రభుత్వ గూఢచర్యం, డేటా చోరీ వల్ల అమెరికా భవిష్యత్తుకు ఎప్పుడూ లేనంత దీర్ఘకాలిక ముప్పు ఉంది” అని అన్నారు...}
\textit{[Output truncated for brevity]} \\
\midrule

TF-IDF V2 &
{\telugufont ఎఫ్‌బీఐ డైరెక్టర్ క్రిస్టఫర్ రే “చైనా ప్రభుత్వ గూఢచర్యం, డేటా చోరీ వల్ల అమెరికా భవిష్యత్తుకు ఎప్పుడూ లేనంత దీర్ఘకాలిక ముప్పు ఉంది” అని అన్నారు...}
\textit{[Output truncated for brevity]} \\
\midrule

mT5 Base &
{\telugufont చైనా ఆర్థిక గూఢచర్యం, డేటా చోరీ వల్ల అమెరికా భవిష్యత్తుకు ఎప్పుడూ లేనంత దీర్ఘకాలిక ముప్పు ఉందని ఫైనాన్షియల్ బ్యాంక్ ఆఫ్ ఇండియా (ఎఫ్‌బీఐ) డైరెక్టర్ చెప్పారు.} \\
\midrule

mT5 Fine-tuned (8-bit LoRA) &
{\telugufont చైనా ప్రభుత్వం గూఢచర్యం, డేటా చోరీ వల్ల అమెరికా భవిష్యత్తుకు ఎప్పుడూ లేనంత దీర్ఘకాలిక ముప్పు ఉందని ఎఫ్‌బీఐ డైరెక్టర్ తెలిపారు.} \\
\bottomrule
\end{tabular}
\end{table}

Three qualitative patterns are visible. First, TF-IDF~V1
and TF-IDF~V2 both return a near-verbatim first sentence
of the article, reflecting the lead-bias position weighting
in V2 and the high term frequency of salient
government-related tokens in both methods. Their stronger
ROUGE-2 performance in aggregate evaluation
(Table~\ref{tab:main}) therefore arises largely from lexical
bigram copying rather than true abstractive compression.

Second, both mT5 variants produce abstractive paraphrases
that deviate lexically from the source and reference,
helping explain their lower ROUGE-2 despite higher
BERTScore. Notably, the mT5~Base output expands the
abbreviation \textit{FBI} as
\textit{{\telugufont ఫైనాన్షియల్ బ్యాంక్ ఆఫ్ ఇండియా}}
(Financial Bank of India), a named-entity hallucination
absent from the source. This illustrates a common
hallucination failure mode in multilingual abstractive
summarization for abbreviation-dense news text.

Third, the mT5~Fine-tuned (8-bit LoRA) output more closely
mirrors the noun-phrase structure of the reference summary,
consistent with the marginal ROUGE-L gain (+0.0032) without
meaningful BERTScore improvement. This suggests that local
fine-tuning primarily refines surface phrasing rather than
semantic content.

\FloatBarrier
%% ============================================================
\section{Limitations and Future Work}
%% ============================================================

This study has several limitations.
(1)~\textbf{No human evaluation:} all evaluations rely on
automatic metrics without human assessment of fluency,
faithfulness, or factual consistency;
(2)~\textbf{No component ablation:} the individual
contributions of character $n$-gram analysis, position
weighting, and adaptive sentence selection in TF-IDF~V2
were not isolated;
(3)~\textbf{TTS quality:} speech synthesis quality was
assessed only qualitatively, without a controlled Mean
Opinion Score (MOS) study;
(4)~\textbf{In-distribution fine-tuning:} LoRA adaptation
was restricted to the news domain, where base model
representations were already strong; and
(5)~\textbf{Domain restriction:} generalizability to other
Telugu text genres remains unverified.

Future work includes out-of-distribution LoRA fine-tuning
(e.g., legal, educational, or conversational Telugu),
component-wise ablation of TF-IDF~V2, longer-context
architectures for documents exceeding 512 tokens, human
correlation studies for metric validation, and extension
to other Dravidian languages.

\section{Conclusion}
%% ============================================================

This paper presented \textit{Saaram}, a deployed Telugu news
summarization system with optional speech synthesis, along
with a comparative study on the complete 1,302-sample
XL-Sum Telugu test split. Our findings show that:
(1)~morphology-aware TF-IDF~V2 yields a 26.6\% ROUGE-2
improvement over TF-IDF~V1, validating character
$n$-grams for agglutinative languages;
(2)~in-distribution 8-bit LoRA fine-tuning provides
negligible semantic gains over the pre-trained mT5 base
model;
(3)~BERTScore provides a more semantically informative
evaluation signal than ROUGE for inflected Telugu outputs,
although human validation remains necessary; and
(4)~the resource-aware router preserves 99.56\% of mT5
quality while providing a reliable fallback under
memory-constrained deployment settings.

%% ============================================================
%%  REFERENCES
%% ============================================================
\begin{thebibliography}{00}

\bibitem{joshi2020text}
A.~Joshi et al., ``Text Summarization for Indian Languages:
A Survey,'' \textit{ACM Trans.\ Asian Low-Resource Lang.\
Inf.\ Process.}, vol.~19, no.~6, pp.~1--35, 2020.

\bibitem{singh2020text}
P.~Singh and K.~Reddy, ``Text Summarization and Document
Summarization using NLP,'' in \textit{Proc.\ ICCIDS},
vol.~167, pp.~745--752, 2020.

\bibitem{hasan2021xlsum}
T.~Hasan et al., ``XL-Sum: Large-Scale Multilingual
Abstractive Summarization for 44 Languages,''
in \textit{Findings ACL-IJCNLP}, pp.~4693--4703, 2021.

\bibitem{srinivasa2019neural}
M.~Srinivasa~Rao and K.~Sridevi, ``Neural Abstractive
Text Summarizer for Telugu Language,''
\textit{IJITEE}, vol.~8, no.~11, pp.~3214--3218, 2019.

\bibitem{anandbabu2023extractive}
G.~L.~Anand Babu and S.~Badugu, ``Extractive Summarization
of Telugu Text Using Modified TextRank and MMR,''
\textit{IJISAE}, vol.~11, no.~2, pp.~512--519, 2023.

\bibitem{erkan2004lexrank}
G.~Erkan and D.~R.~Radev, ``LexRank: Graph-based Lexical
Centrality as Salience in Text Summarization,''
\textit{J.\ Artif.\ Intell.\ Res.}, vol.~22,
pp.~457--479, 2004.

\bibitem{mihalcea2004textrank}
R.~Mihalcea and P.~Tarau, ``TextRank: Bringing Order
into Text,'' in \textit{Proc.\ EMNLP}, pp.~404--411, 2004.

\bibitem{carbonell1998use}
J.~Carbonell and J.~Goldstein, ``The Use of MMR,
Diversity-Based Reranking for Reordering Documents
and Producing Summaries,'' in \textit{Proc.\ ACM SIGIR},
pp.~335--336, 1998.

\bibitem{rush2015neural}
A.~M.~Rush, S.~Chopra, and J.~Weston, ``A Neural Attention
Model for Abstractive Sentence Summarization,''
in \textit{Proc.\ EMNLP}, pp.~379--389, 2015.

\bibitem{nallapati2016abstractive}
R.~Nallapati et al., ``Abstractive Text Summarization
using Sequence-to-Sequence RNNs and Beyond,''
in \textit{Proc.\ CoNLL}, pp.~280--290, 2016.

\bibitem{see2017get}
A.~See, P.~J.~Liu, and C.~D.~Manning, ``Get To The Point:
Summarization with Pointer-Generator Networks,''
in \textit{Proc.\ ACL}, pp.~1073--1083, 2017.

\bibitem{vaswani2017attention}
A.~Vaswani et al., ``Attention is All You Need,''
in \textit{Proc.\ NeurIPS}, pp.~5998--6008, 2017.

\bibitem{raffel2020exploring}
C.~Raffel et al., ``Exploring the Limits of Transfer
Learning with a Unified Text-to-Text Transformer,''
\textit{JMLR}, vol.~21, no.~140, pp.~1--67, 2020.

\bibitem{xue2021mt5}
L.~Xue et al., ``mT5: A Massively Multilingual Pre-trained
Text-to-Text Transformer,'' in \textit{Proc.\ NAACL-HLT},
pp.~483--498, 2021.

\bibitem{kumar2022indicbart}
V.~Kumar et al., ``IndicBART: A Pre-trained Model for
Indic Natural Language Generation,''
in \textit{Findings ACL}, pp.~1849--1863, 2022.

\bibitem{doddapaneni2021primer}
S.~Doddapaneni et al., ``A Primer on Pretrained
Multilingual Language Models,''
arXiv preprint arXiv:2107.00676, 2021.

\bibitem{lin2004rouge}
C.-Y.~Lin, ``ROUGE: A Package for Automatic Evaluation
of Summaries,'' in \textit{Text Summarization Branches Out},
pp.~74--81, 2004.

\bibitem{zhang2020bertscore}
T.~Zhang et al., ``BERTScore: Evaluating Text Generation
with BERT,'' in \textit{Proc.\ ICLR}, 2020.

\bibitem{urlana2022tesum}
A.~Urlana et al., ``TeSum: Human-Generated Abstractive
Summarization Corpus for Telugu,''
in \textit{Proc.\ 13th LREC}, pp.~3862--3869, 2022.

\bibitem{hu2022lora}
E.~J.~Hu et al., ``LoRA: Low-Rank Adaptation of Large
Language Models,'' in \textit{Proc.\ ICLR}, 2022.

\bibitem{dettmers2022int8}
T.~Dettmers et al., ``LLM.int8(): 8-bit Matrix Multiplication
for Transformers at Scale,'' in \textit{Proc.\ NeurIPS}, 2022.

\end{thebibliography}

\end{document}

